\documentclass[12pt,a4paper]{leaflet}
\usepackage[utf8]{inputenc}
\usepackage[german]{babel}
\usepackage{microtype}
\usepackage{graphicx}
\usepackage{color}
\definecolor{gray}{gray}{.5}
\definecolor{darkgray}{gray}{.2}
\setlength{\parindent}{0pt} 
\usepackage{blindtext}
\usepackage[condensed]{tgheros}
\setmargins{2cm}{1cm}{8mm}{8mm}
\usepackage[absolute]{textpos}

\usepackage{hyperref}
\hypersetup{
    pdftoolbar=true,
    pdfmenubar=true,
    pdffitwindow=false,
    pdfstartview={FitH},
    pdfnewwindow=true,
    colorlinks=true,
    linkcolor=black,
    citecolor=black,
    filecolor=black,
    urlcolor=black
}
\def\UrlFont{\bfseries}
\pagestyle{empty} 
\AddToBackground*{2}{\put(0,25){\includegraphics[scale=1]{welle2}}}

\newcommand{\software}[7]{
\begin{minipage}[t][0.4\textheight]{1\textwidth}
\flushleft
\subsection*{#1 \hfill \textcolor{gray}{#2}}
#3 
\medskip

%\scalebox{.9}[1.0]{\textit{Betriebssystem:} \textbf{#4}} \\
\textit{Betriebssystem:} \textbf{#4} \\
\textit{Lizenz:} \textbf{#6} \\
\textit{Deutschsprachig:} \textbf{#7} \\
\textit{Homepage:} \url{#5} \\
\vglue  2 true cm
\end{minipage}
}

% Hack for LibreOffice/OpenOffice: Version of \software without link
\newcommand{\softwaretwo}[7]{
\begin{minipage}[t][0.4\textheight]{1\textwidth}
\flushleft
\subsection*{#1 \hfill \textcolor{gray}{#2}}
#3 
\medskip

%\scalebox{.9}[1.0]{\textit{Betriebssystem:} \textbf{#4}} \\
\textit{Betriebssystem:} \textbf{#4} \\
\textit{Lizenz:} \textbf{#6} \\
\textit{Deutschsprachig:} \textbf{#7} \\
\textit{Homepage:} \textbf{#5} \\
\vglue  2 true cm
\end{minipage}
}

\newcommand{\frontpage}[1]{

\begin{titlepage}
\begin{flushleft}
\vglue  2 true cm

\textsc{\Huge Freie Software} 
\newcommand{\HRule}{\rule{0.9\linewidth}{0.5mm}}
\HRule \\[0.4cm]
\textsc{\textbf{\Huge #1}}

\vglue  2 true cm

{ \large\it
Acht Freie \\
Computerprogramme, \\
die Deine Rechte \\
als Nutzer/-in \\
respektieren \\
und schützen. }
\end{flushleft}
\vfill
\begin{flushright}
\begin{minipage}{3.5cm}
\center
\includegraphics[scale=0.15]{plussy} \\
\bf FSFE-Fellowship\\ München
\end{minipage}
\end{flushright}
\end{titlepage}

\setlength{\TPHorizModule}{\textwidth}
\setlength{\TPVertModule}{\textheight}
\begin{textblock}{1}(0.1,1)
\small\textcolor{darkgray}{
Impressum: \\
FSFE Fellowship München \\
c/o Free Software Foundation Europe e.V. \\
Schönhauser Allee 6/7 \\
10119 Berlin}
\end{textblock}

\flushleft
}

\newcommand{\backpage}{
\newpage
\subsection{Was ist »Freie Software«?}
Von »Freier Software« spricht man, sobald ein Software-Hersteller
sein \textbf{Computerprogramm} unter einer \textbf{Freien Softwarelizenz} veröffentlicht.

Im Gegensatz zu proprietären Softwarelizenzen
(welche Art und Weise der Softwarenutzung meist
zugunsten des Herstellers einschränken), \textbf{zielen
Freie Softwarelizenzen} darauf ab, \textbf{Deine Rechte
als Nutzer/-in} zu verteidigen. Diese Rechte können
anhand von \textbf{vier Freiheiten} beschrieben werden:

\begin{itemize}
\item Der Freiheit, dass Programm jederzeit und für jeden beliebigen Zweck zu \textbf{verwenden}
\item Der Freiheit, die Funktionsweise des Programms zu \textbf{verstehen}, indem Du seinen Programmcode untersuchst
\item Der Freiheit, das Programm zu \textbf{verändern}, um es Deinen Anforderungen und Wünschen anzupassen
\item Der Freiheit, Kopien des (un)veränderten Programms an Deine Mitmenschen zu \textbf{verteilen}, damit auch andere davon profitieren können
\end{itemize}

Weitere Informationen zu Freier Software,
Freien Softwarelizenzen und dem »Free Software Movement«
findest Du unter \url{www.fsfe.org}.

}
